\DocumentMetadata{
  pdfversion=2.0,
  pdfstandard=ua-2,
  lang=en-US,
  tagging=on,
  tagging-setup={math/setup=mathml-SE,tabsorder=structure}
}
\documentclass[11pt,letterpaper]{article}
\usepackage[margin=0.68in,includefoot]{geometry}
\usepackage{newcomputermodern}
\usepackage{amsmath,amscd,mathtools}
\usepackage{tikz,adjustbox}
\providecommand{\square}{\mdlgwhtsquare}
\usepackage[hidelinks]{hyperref}
\hypersetup{
  pdftitle={Analysis PhD Exam, May 2005},
  pdfauthor={University of Florida Department of Mathematics},
  pdfsubject={Graduate mathematics examination},
  pdfdisplaydoctitle=true,
  bookmarksopen=true
}
\setcounter{secnumdepth}{0}
\setlength{\parindent}{0pt}
\setlength{\parskip}{0.28em}
\setlength{\emergencystretch}{3em}
\setlength{\leftmargini}{2.15em}
\setlength{\leftmarginii}{2.65em}
\setlength{\leftmarginiii}{3em}
\pagestyle{empty}
\raggedbottom
\allowdisplaybreaks
\NewTaggingSocketPlug{block/list/label}{examproblem}{%
  \tagstructbegin{tag=Lbl}\tagstructend%
  \tagstructbegin{tag=\LBody}%
  \tagstructbegin{tag=\ProblemHeadingTag,title-o={\ProblemHeadingText}}%
  \tagmcbegin{tag=\ProblemHeadingTag,actualtext={\ProblemHeadingText}}%
  #2%
  \tagmcend%
  \tagstructend%
}
\newenvironment{examproblems}[2]{%
  \def\ProblemHeadingTag{#2}%
  \AssignTaggingSocketPlug{block/list/label}{examproblem}%
  \begin{enumerate}%
  \setcounter{enumi}{#1}%
  \renewcommand{\labelenumi}{\textbf{\theenumi.}}%
  \setlength{\topsep}{0.35em}%
  \setlength{\partopsep}{0pt}%
  \setlength{\itemsep}{0.48em}%
  \setlength{\parsep}{0.18em}%
}{\end{enumerate}}
\newenvironment{examsubparts}{%
  \AssignTaggingSocketPlug{block/list/label}{default}%
  \begin{enumerate}%
  \setlength{\topsep}{0.18em}%
  \setlength{\partopsep}{0pt}%
  \setlength{\itemsep}{0.16em}%
  \setlength{\parsep}{0.1em}%
}{\end{enumerate}}
\newenvironment{examinstructions}{%
  \par\begingroup\itshape
}{%
  \par\endgroup\vspace{0.18em}
}
\makeatletter
\renewcommand\subsection{\@startsection{subsection}{2}{0pt}%
  {-1.25ex plus -.25ex minus -.1ex}%
  {0.55ex plus .1ex}%
  {\normalfont\large\bfseries}}
\renewcommand{\@maketitle}{%
  \newpage\null\vskip -1em%
  {\footnotesize\bfseries
    \href{https://www.ufl.edu/}{University of Florida}, \href{https://math.ufl.edu/}{Department of Mathematics}\par}%
  \vskip -0.5em%
  \tagstructbegin{tag=H1,title={Analysis PhD Exam, May 2005}}%
  \tagpdfparaOff%
  \tagmcbegin{tag=H1}%
    {\LARGE\bfseries\href{https://gma.math.ufl.edu/exams/analysis-phd/}{Analysis PhD Exam}, May 2005\par}%
  \tagmcend%
  \tagpdfparaOn%
  \tagstructend%
  \vskip 0.25em%
  \tagmcbegin{artifact}%
    \hrule height 1pt%
  \tagmcend%
  \vskip 1em%
}
\makeatother
\title{Analysis PhD Exam, May 2005}
\author{}
\date{}
\begin{document}
\maketitle
\thispagestyle{empty}
\begin{examinstructions}
Give complete proofs and computations. Be particularly careful to indicate why the most crucial steps in your proof are true.
\end{examinstructions}
\begin{examproblems}{0}{H2}
\def\ProblemHeadingText{Problem 1.}
\item
Let~\(\mu\) be a positive Radon measure on a locally compact Hausdorff space~\(T\). Let \(f_n:T\to\overline{\mathbb R}\) be~\(\mu\)-measurable for all \(n\in\mathbb N\). Prove that \(\sup\{f_n\mid n\in\mathbb N\}\) is~\(\mu\)-measurable.
\def\ProblemHeadingText{Problem 2.}
\item
Let~\(E\) be a complex Banach space. Suppose \(\{x_n\}_{n\in\mathbb N}\subset E\) converges in the weakened topology \(\sigma(E,E')\) to \(x\in E\), i.e. \(\lvert\langle x_n-x,y\rangle\rvert\to0\) for all \(y\in E'\) as \(n\to\infty\). Prove that~\(x\) is contained in the closure of the linear span of the set \(\{x_n\}_{n\in\mathbb N}\).
\def\ProblemHeadingText{Problem 3.}
\item
Let~\(E\) be a topological vector space. Prove the following:
\begin{examsubparts}
\item[\textnormal{(a)}]
If \(A\subset E\) is arbitrary and \(U\subset E\) is open, then \(A+U\) is open.
\item[\textnormal{(b)}]
If \(A\subset E\) is compact and \(B\subset E\) is closed, then \(A+B\) is closed.
\item[\textnormal{(c)}]
The sum of closed subsets may fail to be closed.
\end{examsubparts}
\def\ProblemHeadingText{Problem 4.}
\item
Let \(\delta\in\mathcal D'(\mathbb R)\) denote the distribution given by \(\delta(\phi)=\phi(0)\), for all \(\phi\in\mathcal D(\mathbb R)\). Prove your answers are correct.
\begin{examsubparts}
\item[\textnormal{(a)}]
For which \(f\in C^\infty(\mathbb R)\) is \(f\cdot\delta'=0\)?
\item[\textnormal{(b)}]
What is the support of~\(\delta\)?
\item[\textnormal{(c)}]
If \(X\in\mathcal D'(\mathbb R)\), \(f\in C^\infty(\mathbb R)\), and~\(f\) restricted to the support of~\(X\) is~\(0\), is it true that \(f\cdot X=0\)?
\end{examsubparts}
\def\ProblemHeadingText{Problem 5.}
\item
Let \(K:[0,1]\times[0,1]\to\mathbb R\) be continuous and \(f:[0,1]\times\mathbb R\to\mathbb R\) be continuous and bounded. Show that the Hammerstein equation
\[
u(s)=\int_0^1 K(s,t)f(t,u(t))\,dt,\qquad s\in[0,1],
\]
 has a solution in \(E=\{x:[0,1]\to\mathbb R\mid x\text{ continuous}\}\).
\def\ProblemHeadingText{Problem 6.}
\item
Let \(1\in\mathcal D'(\mathbb R)\) be given by \(1(\phi)=\int_{\mathbb R}\phi\,dx\), for all \(\phi\in\mathcal D(\mathbb R)\). Compute \((1*\delta')(\phi)\) for all \(\phi\in\mathcal D(\mathbb R)\).
\def\ProblemHeadingText{Problem 7.}
\item
Let~\(\mu\) be a~\(\sigma\)-finite positive Radon measure on a locally compact Hausdorff space~\(T\). Let \(\{f_n\}_{n\in\mathbb N}\subset L^p(T,\mu)\) for some \(p\in[1,\infty)\). Suppose that \(\lim_{n\to\infty}f_n(t)=f(t)\) \(\mu\)-a.e., that \(f\in L^p(T,\mu)\), and that \(\lim_{n\to\infty}\int\lvert f_n\rvert^p\,d\mu=\int\lvert f\rvert^p\,d\mu\). Prove that \(\lim_{n\to\infty}\int\lvert f_n-f\rvert^p\,d\mu=0\). Hint: Bring the Lebesgue Dominated Convergence Theorem and the~\(\sigma\)-finiteness of the measure into play.
\def\ProblemHeadingText{Problem 8.}
\item
Let \(\mathcal H\) be a Hilbert space with inner product \((\cdot,\cdot)\), and let \(\{y_\alpha\}_{\alpha\in A}\subset\mathcal H\) be an orthonormal family. Let \(x\in\mathcal H\).
\begin{examsubparts}
\item[\textnormal{(a)}]
Prove that \(\sum_\alpha (x,y_\alpha)y_\alpha\) converges strongly to an element \(y\in\mathcal H\).
\item[\textnormal{(b)}]
Prove that the vector \(x-y\) is orthogonal to the closure of the linear span of the set \(\{y_\alpha\}_{\alpha\in A}\).
\end{examsubparts}
\end{examproblems}
\end{document}
